%%
%% This is file `sample-manuscript.tex',
%% generated with the docstrip utility.
%%
%% The original source files were:
%%
%% samples.dtx  (with options: `manuscript')
%% 
%% IMPORTANT NOTICE:
%% 
%% For the copyright see the source file.
%% 
%% Any modified versions of this file must be renamed
%% with new filenames distinct from sample-manuscript.tex.
%% 
%% For distribution of the original source see the terms
%% for copying and modification in the file samples.dtx.
%% 
%% This generated file may be distributed as long as the
%% original source files, as listed above, are part of the
%% same distribution. (The sources need not necessarily be
%% in the same archive or directory.)
%%
%% Commands for TeXCount
%TC:macro \cite [option:text,text]
%TC:macro \citep [option:text,text]
%TC:macro \citet [option:text,text]
%TC:envir table 0 1
%TC:envir table* 0 1
%TC:envir tabular [ignore] word
%TC:envir displaymath 0 word
%TC:envir math 0 word
%TC:envir comment 0 0
%%
%%
%% The first command in your LaTeX source must be the \documentclass command.
%\documentclass[sigconf,review,anonymous]{acmart}
%\documentclass[manuscript,review,anonymous]{acmart}
%\documentclass[manuscript,review,anonymous]{acmart}
\documentclass[manuscript,review]{acmart}



% \usepackage[vskip=1em,font=itshape,leftmargin=2em,rightmargin=2em]{quoting}

%%
%% \BibTeX command to typeset BibTeX logo in the docs
% \AtBeginDocument{%
%   \providecommand\BibTeX{{%
%     \normalfont B\kern-0.5em{\scshape i\kern-0.25em b}\kern-0.8em\TeX}}}

%% Rights management information.  This information is sent to you
%% when you complete the rights form.  These commands have SAMPLE
%% values in them; it is your responsibility as an author to replace
%% the commands and values with those provided to you when you
%% complete the rights form.
% \copyrightyear{2024}
% \acmYear{2024}
% \setcopyright{rightsretained}
% \acmConference[Conference acronym 'XX]{Make sure to enter the correct
%   conference title from your rights confirmation emai}{June 03--05,
%   2018}{Woodstock, NY}
% \acmISBN{978-1-4503-XXXX-X/18/06}

%% These commands are for a PROCEEDINGS abstract or paper.
% \acmConference[Conference acronym 'XX]{Make sure to enter the correct
%   conference title from your rights confirmation emai}{June 03--05,
%   2018}{Woodstock, NY}
% \acmISBN{978-1-4503-XXXX-X/18/06}


%%
%% Submission ID.
%% Use this when submitting an article to a sponsored event. You'll
%% receive a unique submission ID from the organizers
%% of the event, and this ID should be used as the parameter to this command.
%%\acmSubmissionID{123-A56-BU3}

%%
%% For managing citations, it is recommended to use bibliography
%% files in BibTeX format.
%%
%% You can then either use BibTeX with the ACM-Reference-Format style,
%% or BibLaTeX with the acmnumeric or acmauthoryear sytles, that include
%% support for advanced citation of software artifact from the
%% biblatex-software package, also separately available on CTAN.
%%
%% Look at the sample-*-biblatex.tex files for templates showcasing
%% the biblatex styles.
%%

%%
%% The majority of ACM publications use numbered citations and
%% references.  The command \citestyle{authoryear} switches to the
%% "author year" style.
%%
%% If you are preparing content for an event
%% sponsored by ACM SIGGRAPH, you must use the "author year" style of
%% citations and references.
%% Uncommenting
%% the next command will enable that style.
%%\citestyle{acmauthoryear}

%%
%% end of the preamble, the start of the body of the document source.



\begin{document}

%%
%% The "title" command has an optional parameter,
%% allowing the author to define a "short title" to be used in page headers.
% \title{Mixed Reality Street Play: Towards a Playful City with Ubiquitous Mixed Reality Head-Mounted Display}

\title{"Don't Die!": Performing Cybroc}

% Synchronous Asymmetric

%%
%% The "author" command and its associated commands are used to define
%% the authors and their affiliations.
%% Of note is the shared affiliation of the first two authors, and the
%% "author note" and "authornotemark" commands
%% used to denote shared contribution to the research.

\author{Ke Huang}
\orcid{0009-0000-6884-9076}
\affiliation{%    
    \institution{China Academy of Art}
    \city{Hangzhou}
    \country{China}
}
\affiliation{
    \institution{Tsinghua University}
    \city{Beijing}
    \country{China}
}
\email{kehuang81@gmail.com}

\author{Botao Amber Hu}\authornote{Corresponding author}
\orcid{0000-0002-4504-0941}
\affiliation{%
  \institution{Reality Design Lab}
  \city{New York City}
  \country{USA}
  }
\affiliation{%
  \institution{University of Oxford}
  \city{Oxford}
  \country{UK}
  }
\email{botao@reality.design}


%%
%% By default, the full list of authors will be used in the page
%% headers. Often, this list is too long, and will overlap
%% other information printed in the page headers. This command allows
%% the author to define a more concise list
%% of authors' names for this purpose.
% \renewcommand{\shortauthors}{Botao Amber Hu et al}

%%
%% The abstract is a short summary of the work to be presented in the
%% article.
\begin{abstract}

\end{abstract}

%% The code below is generated by the tool at http://dl.acm.org/ccs.cfm.
%% Please copy and paste the code instead of the example below.
%%
\begin{CCSXML}
<ccs2012>
<concept>
  <concept_id>10003120.10003123.10011758</concept_id>
  <concept_desc>Human-centered computing~Interaction design theory, concepts and paradigms</concept_desc>
  <concept_significance>100</concept_significance>
</concept>
<concept>
  <concept_id>10003120.10003121.10003124.10010392</concept_id>
  <concept_desc>Human-centered computing~Mixed / augmented reality</concept_desc>
  <concept_significance>500</concept_significance>
</concept>
<concept>
<concept_id>10003120.10003130.10003233</concept_id>
<concept_desc>Human-centered computing~Collaborative and social computing systems and tools</concept_desc>
<concept_significance>100</concept_significance>
</concept>
</ccs2012>
\end{CCSXML}

\ccsdesc[100]{Human-centered computing~Interaction design theory, concepts and paradigms}
\ccsdesc[500]{Human-centered computing~Mixed / augmented reality}
\ccsdesc[100]{Human-centered computing~Collaborative and social computing systems and tools}

%%
%% Keywords. The author(s) should pick words that accurately describe
%% the work being presented. Separate the keywords with commas.
\keywords{}

% \received{21 February 2024}
% \received[revised]{21 February 2024}
% \received[accepted]{21 February 2024}

\begin{teaserfigure}
    \centering
    \Description{}
    \label{fig:teaser}
\end{teaserfigure}

%%
%% This command processes the author and affiliation and title
%% information and builds the first part of the formatted document.
\maketitle


\section{Introduction}

\section{Background}

\subsection{Longevity Movement and Transhumanism}

<Prompt>: talk about Bryan Johnson, Longevity Movement
<Prompt>: big background of Transhumanism


\subsection{Performing Economy and Memetics}

\subsubsection{Meme spreading and Capitalism}
<Prompt>: talk about How Bryan Johnson Utilize the Meme spreading to make it a movement 

\subsubsection{Live Streaming and Memecoin}
<Prompt>: talk about How Bryan Johnson Utilize the Meme spreading to make it a movement 



\subsection{Cybroc: }
What's Cybroc?

\section{Performing Cybroc}
The Process 

\section{Discussion}

\subsection{Performing Satire}
\subsection{Nature vs Artificial Life}
\subsection{Ethical Dilemma}


\section{Conclusion}








%%
%% The next two lines define the bibliography style to be used, and
%% the bibliography file.
\bibliographystyle{ACM-Reference-Format}
\bibliography{reference}

%%
%% If your work has an appendix, this is the place to put it.
\appendix
%\input{09appendix.tex}


\end{document}
\endinput
%%
%% End of file `sample-manuscript.tex'.
